\chapter{Introduction}

\section{Context and Motivation}\label{sec:context}
Participatory democracy platforms such as Decidim~\cite{barandiaran2024decidim}, Cap
Collectif~\cite{shin2024digitaltools}, and Consul~\cite{aranacatania2021consul} enable thousands
of citizens to propose and debate public policy. Their adoption has been substantial: Decide
Madrid, built on Consul, saw over 430{,}000 users submit more than 26{,}000 proposals and
125{,}000 comments in its citizen-proposals process alone~\cite{aranacatania2021consul}. Similar
platforms are now deployed by city and national governments across Europe and Latin America for
participatory budgeting, collaborative law-making, and large-scale public consultations such as
France's Grand D\'ebat National.

Their success creates a problem of scale. As the number of contributions grows, no citizen can
read the full debate, so they cannot tell which existing proposals resemble or oppose their own,
and no organizer can manually review thousands of comments to determine where consensus or
disagreement lies. This is precisely the difficulty documented on Consul-based platforms, where
both citizens and administrators reported being unable to obtain an overview of the volume of
proposals and comments generated~\cite{aranacatania2021consul}. The result is a form of
information overload that risks undermining the very inclusiveness these platforms are meant to
provide, since a citizen's contribution is only meaningfully heard if it can be situated within
the broader debate, and a debate too large to read is, in practice, a debate too large to
participate in.

Prior work has shown that some of this structure can be recovered automatically. Natural language
processing models can detect argumentative relations between citizen proposals, determining
whether two proposals make the same point or whether they are in direct conflict, and these
relations can be organized into an argumentative graph that makes the shape of the debate explicit
rather than leaving it buried in unstructured text. Helping a user see this shape, meaning the
relevant arguments on both sides of an issue without redundancy, is fundamentally a retrieval
problem: given a policy under debate, which arguments should be surfaced, and how should they be
organized so that the retrieval itself does not reintroduce the very redundancy and imbalance it
is meant to resolve?

\section{Problem Statement}
\label{sec:problem}

The preceding section motivates the need to make the structure of a large-scale civic debate
explicit rather than leaving it as an unstructured collection of proposals. The remaining
question is how a retrieval system should exploit that structure once it exists.

The default approach, dense retrieval over text embeddings, treats each argument as an
independent point in a vector space and retrieves the arguments closest to a query. This is
effective at finding topical relevance, but it has no notion of two arguments being redundant
restatements of the same point, nor of two arguments being in direct tension. A user asking to
see the debate on a policy is not well served by five arguments that all reflect the same view
in different words, nor by a set of pro and con arguments that never actually engage the same
underlying claim. Recovering this structure, which arguments repeat each other and which
arguments conflict, requires relations between arguments, not just positions in an embedding
space.

Graph-based retrieval methods (Graph-RAG) address this by building an explicit graph over the
corpus and using its structure, rather than embedding distance alone, to guide retrieval. Whether
this actually improves retrieval for argumentative content is, however, an open empirical
question rather than an established result: the graph adds a strong inductive bias (that
structurally connected arguments are the ones worth retrieving together), and inductive biases
help only when they match the true structure of the task. This thesis investigates that question
directly. Concretely, it builds an argument graph over a large debate corpus, with edges
representing contradiction and equivalence between arguments, and compares retrieval that
exploits this graph structure against retrieval that relies on embedding similarity alone, across
two tasks: retrieving a balanced, non-redundant, and genuinely conflicting set of arguments for a
policy (conflict-map retrieval), and producing a debate summary from the retrieved arguments
(debate summarization).

\section{Research Questions}
\label{sec:rq}

The problem statement above is investigated through three research questions, each targeting a
distinct threat to the validity of a simple positive answer.

\begin{description}
\item[RQ1 (Does structure help).] For conflict-map retrieval (UC1), does a graph-structure-aware
  retrieval system, one that expands its candidate pool through explicit \textsc{contradicts}
  relations and ranks by graph centrality, achieve higher relevance and diversity than a system
  that relies on embedding similarity alone, and does any such advantage survive into debate
  summarization (UC2)?

\item[RQ2 (Is the advantage genuine).] A relevance or diversity gap measured on a corpus's native
  pool can reflect genuine retrieval quality, or it can reflect properties of the evaluation setup
  itself, such as which candidates happen to populate the pool and how they were selected. Under
  an evaluation protocol deliberately hardened with adversarial distractors, does any advantage
  observed under RQ1 persist, and is it attributable to the retrieval method rather than to how
  the pool was constructed?

\item[RQ3 (Does the finding generalize).] The corpus used to answer RQ1 and RQ2 is a single
  English-language, crowd-authored dataset. Does the same pattern of results hold on an
  independently constructed corpus of real citizen contributions in a different language, namely
  the French Grand D\'ebat National, or is it specific to the properties of the original corpus?
\end{description}

RQ1 is answered in Chapter~\ref{ch:results} through the UC1 and UC2 evaluations. RQ2 is answered
by two robustness checks built into the same evaluation protocol: injecting distractor arguments
with their own graph edges, and reconstructing the candidate pool under random rather than
similarity-based selection. RQ3 is answered by repeating the UC1 pipeline on the Grand D\'ebat
National corpus described in Chapter~\ref{ch:corpus}. Taken together, these three questions are
designed so that a positive answer to RQ1 alone would not be sufficient to claim that graph
structure improves retrieval: the claim is only supported if it also survives RQ2 and, ideally,
RQ3.

\section{Technical Background}
\label{sec:background}

This thesis draws on four strands of prior work, each contributing one layer of the system built
here.

\subsection{Natural language processing and argumentative relations.}
\label{sec:background-nlp}

Determining whether one piece of text supports, restates, or contradicts another is a core NLP
problem, most commonly framed as natural language inference (NLI): given a premise and a
hypothesis, predict whether the hypothesis is entailed, contradicted, or neutral with respect to
the premise. Argument mining specializes this general machinery to argumentative text
specifically, extracting claims, premises, and the relations between them from unstructured
documents~\cite{lippi2016argumentation}. Prior work has shown that such relations, including
direct attack and support between arguments, can be identified automatically using supervised
learning over textual features~\cite{cocarascu2017}, and that entailment-style reasoning
transfers naturally to the argumentative setting, where two arguments can be treated as
equivalent, contradictory, or unrelated depending on their entailment relationship in each
direction~\cite{cabrio2012}. This thesis relies on the same underlying idea, using NLI as the
scoring mechanism that produces the \textsc{contradicts} and \textsc{equivalent} relations
described in Chapter~\ref{ch:methodology}.

\subsection{Knowledge graphs.}
A knowledge graph represents entities as nodes and their relations as typed edges, making
relational structure explicit and queryable rather than leaving it implicit in text or buried in
independent embeddings~\cite{hogan2021knowledgegraphs}. The corpus used in this thesis is
organized exactly this way: arguments are nodes, and \textsc{contradicts}/\textsc{equivalent}
relations are typed edges between them, forming an argumentative knowledge graph. Graph
structure of this kind supports operations that are awkward or impossible with a flat set of
embeddings alone, such as measuring how central an argument is within the debate through
graph centrality, or finding the community of arguments that all restate the same underlying
point through community detection.

\subsection{Retrieval-augmented generation.}
Retrieval-augmented generation (RAG) combines a retriever, which selects relevant items from a
corpus given a query, with a generator, typically a large language model, that conditions its
output on the retrieved items~\cite{lewis2020rag}. RAG was originally motivated by knowledge
grounding, letting a generator produce answers backed by retrieved evidence rather than relying
solely on parametric memory, and standard RAG retrieves items independently by embedding
similarity: each candidate is scored against the query in isolation, with no notion of how
candidates relate to one another.

\subsection{Graph-RAG.}
Graph-RAG extends this paradigm by making the retriever graph-aware: rather than treating
candidates as independent points in an embedding space, it exploits explicit relations between
them, using graph traversal, centrality, or community structure to decide what to retrieve and how
to diversify it~\cite{edge2024graphrag}. This matters directly for the retrieval problem posed in
Section~\ref{sec:problem}: standard RAG can retrieve five arguments that are all close to a query
in embedding space without ever detecting that they restate the same point, since embedding
similarity alone does not distinguish topical closeness from redundancy or conflict. A
graph-aware retriever can use \textsc{equivalent} edges to avoid retrieving restatements of the
same argument, and \textsc{contradicts} edges to actively seek out arguments in genuine tension,
which is precisely the conflict-map retrieval task (UC1) this thesis evaluates. Whether this
theoretical advantage translates into a measurable empirical improvement over standard RAG, on a
real argumentative corpus rather than the general-domain Wikipedia-style corpora typically used
to evaluate Graph-RAG, is the question this thesis investigates.


\subsection{A Brief History: Why These Techniques Became Prominent}
\label{sec:background-history}

RAG and Graph-RAG are recent enough that understanding why they became prominent, rather than
only what they do, is useful context for the design choices made later in this thesis.

The RAG architecture itself was proposed in 2020 as an academic technique for knowledge-intensive
question answering~\cite{lewis2020rag}, but it remained a comparatively specialized topic for
roughly three years. Its move from a research technique to a widely adopted engineering practice
followed the public release of general-purpose conversational LLMs in late 2022, whose fluency
made a specific weakness visible at scale: these models generate confident, well-formed answers
that can nonetheless be factually wrong, a failure mode generally termed hallucination. A model's
parametric knowledge is also frozen at training time and cannot easily incorporate new or
private information without retraining. RAG addresses both problems at once: conditioning
generation on retrieved passages gives the model something to ground its answer in, and updating
the corpus updates the model's effective knowledge without touching its parameters at all. By
2023, retrieval had become a standard component of production LLM systems, and RAG in particular
had gone from a specialized technique to one of the most widely deployed patterns for building
LLM applications~\cite{gao2023ragsurvey}.

Graph-RAG emerged as a response to a specific limitation of this now-standard RAG pattern: flat
vector retrieval struggles with questions that require synthesizing information spread across
many documents, rather than answerable from any single retrieved passage, since similarity
search surfaces individually relevant items without any notion of how those items relate to each
other collectively. Microsoft Research's Graph-RAG system, which builds a knowledge graph over a
corpus and summarizes at the level of detected communities to answer such
global, corpus-spanning questions, was an influential proposal in 2024 for addressing this
gap~\cite{edge2024graphrag} and popularized the term ``Graph-RAG'' for the broader family of
graph-structured retrieval approaches that has since grown around it.

For this thesis, the significant part of this history is what it implies about maturity. RAG is
well-established, with a large literature on retriever design, chunking, and evaluation, so
System~A in this thesis builds on settled practice. Graph-RAG is comparatively young, its
empirical evaluations concentrate on general-domain, entity-centric corpora such as
Wikipedia-derived text, and rather less is established about how it performs on argumentative
text, where the relations that matter (contradiction and equivalence between claims) are not the
entity co-occurrence relations that most existing Graph-RAG systems are built around. This is the
specific gap this thesis addresses.

\subsection{How RAG Works}
\label{sec:background-rag-howitworks}

A RAG system has two components. The \emph{retriever} takes a query and returns the $k$ items
from a corpus most relevant to it, typically by embedding the query and every corpus item into a
shared vector space and ranking candidates by cosine similarity~\cite{lewis2020rag}. The
\emph{generator}, usually a large language model, then conditions its output on both the original
query and the retrieved items, rather than relying only on knowledge stored in its parameters.
This separation is the central idea: retrieval supplies up-to-date, specific, and verifiable
content, while generation supplies fluent, task-adapted language production. Because retrieval and
generation are decoupled, the corpus can be updated without retraining the generator, and the
system's outputs can, in principle, be traced back to the specific items retrieved.

A well-known limitation of naive top-$k$ retrieval is redundancy: the $k$ most similar items to a
query are often near-duplicates of each other, since text that is highly similar to the query
tends also to be highly similar to other similarly-relevant text. Maximal marginal relevance
(MMR) addresses this by re-ranking candidates to jointly maximize relevance to the query and
dissimilarity to items already selected~\cite{carbonell1998mmr}, trading off some relevance for
broader coverage of the retrieved set. This diversity/relevance tradeoff, and the metrics used to
quantify it, recur throughout this thesis's evaluation methodology (Chapter~\ref{ch:evaluation}).

\subsection{How Graph-RAG Works}
\label{sec:background-graphrag-howitworks}

Graph-RAG replaces or augments the flat vector index with an explicit graph over the corpus,
where nodes are corpus items and edges encode a relation between them, such as citation,
co-reference, shared entities, or, in this thesis, argumentative agreement and disagreement.
Retrieval then exploits this structure directly. Two broad strategies are common in the
literature. The first constructs communities of densely connected nodes and, at query time,
retrieves or summarizes at the level of communities rather than individual items, letting the
system answer questions that require synthesizing information spread across many related items
rather than any single one~\cite{edge2024graphrag}. The second, closer to the approach used in
this thesis, expands a candidate set by traversing edges from an initial seed, so that items
connected to a highly relevant node are considered even if they are not themselves close to the
query in embedding space, and ranks the resulting candidates using graph-centrality measures such
as PageRank in addition to, or instead of, embedding similarity.

The practical payoff is that Graph-RAG can retrieve based on \emph{relationships} that embedding
similarity does not capture. Two arguments that directly contradict each other are not
necessarily close in embedding space, since surface wording can differ substantially even when
the underlying claims are in direct tension; conversely, two arguments can be embedding-close
while being logically unrelated, simply because they discuss the same topic. A graph edge encodes
the relationship explicitly, rather than leaving the retriever to infer it indirectly from
distance in a vector space.

\subsection{Advantages and Limitations}
\label{sec:background-tradeoffs}

Both paradigms have genuine strengths and corresponding weaknesses, summarized in
Table~\ref{tab:rag-vs-graphrag}.

Standard RAG is simple to build and maintain: it requires only an embedding model and a vector
index, degrades gracefully as the corpus grows, and needs no relation-extraction step. Its
weakness is exactly what motivates Graph-RAG: it has no mechanism for reasoning about how
candidates relate to one another, so redundancy and conflict must be inferred indirectly, if at
all, from embedding geometry.

Graph-RAG's strength mirrors this weakness: with a well-constructed graph, it can retrieve based
on explicit relational signal that embedding similarity structurally cannot represent. Its cost is
that this strength is entirely conditional on graph quality. Building the graph requires a
relation-extraction step (Section~\ref{sec:background-nlp}) that can introduce its own errors, an
under-connected or noisy graph provides little advantage over standard RAG, and the additional
machinery, edge construction, graph traversal, centrality computation, adds engineering and
computational cost that standard RAG does not incur. Whether this cost is worth paying is an
empirical question that depends on the corpus and the task, not a foregone conclusion in
Graph-RAG's favor, which is precisely the question this thesis investigates for argumentative
retrieval.

\begin{table}[ht]
\centering
\caption{RAG versus Graph-RAG: general properties.}
\label{tab:rag-vs-graphrag}
\begin{tabularx}{\linewidth}{@{}l X X@{}}
\toprule
\textbf{Property} & \textbf{Standard RAG} & \textbf{Graph-RAG} \\
\midrule
Candidate comparison & Each candidate scored against the query independently & Candidates can be scored relative to one another via graph structure \\
Redundancy handling & Indirect, via embedding-distance heuristics (e.g.\ MMR) & Direct, via explicit relation edges (e.g.\ equivalence) \\
Relational reasoning & None; relationships must be inferred from geometry & Explicit; relationships are given as typed edges \\
Setup cost & Low: embedding model and vector index & Higher: requires a relation-extraction pipeline \\
Failure mode & Redundant or topically narrow results & Degrades to no better than standard RAG if the graph is sparse or noisy \\
\bottomrule
\end{tabularx}
\end{table}

\section{Thesis Structure}